% Online Resource 1 for the Sports Engineering submission candidate
\documentclass[pdflatex,sn-mathphys-num]{sn-jnl}

\usepackage{graphicx}
\usepackage{amsmath,amssymb}
\usepackage{booktabs}
\usepackage{multirow}
\usepackage{longtable}

\raggedbottom

\begin{document}

\title{Online Resource 1: Detailed ablations and robustness analyses for
calibration-free two-view 3D pose fusion}

\author*[1]{\fnm{Kaixu} \sur{Chen}}\email{chenkaixusan@gmail.com}
\affil*[1]{\orgname{CCS}}

\maketitle

\section{Purpose}

This Online Resource reports prespecified and post hoc ablations, synthetic
corruption settings, robustness diagnostics and secondary cohort results for
the accompanying article. It preserves the separation between training
evidence and evaluation evidence: triangulated poses were resolved only by the
evaluation layer and were never used for model training, validation or
checkpoint selection.

\section{Ablation definitions}

\begin{table}[h]
\caption{Ablation design. A4--A9 share the same network. Checkmarks indicate
active loss groups. A7--A9 are post hoc trunk-twist variants and do not replace
the prespecified complete model A6.}
\label{tab:ablation}
\centering
\small
\setlength{\tabcolsep}{4pt}
\begin{tabular}{clp{0.46\linewidth}}
\toprule
ID & Method & Active components\\
\midrule
A0 & Face stream only & Deterministic single-view output\\
A1 & Side stream in face reference & Deterministic single-view output\\
A2 & Canonical arithmetic mean & Deterministic mean\\
A3 & Deterministic quality mean & Fixed quality weighting\\
A4 & Spatial-objective fusion & Learned residual; recovery, identity and skeletal losses\\
A5 & Rotation-temporal fusion & A4 plus rotation and temporal losses\\
A6 & Complete self-supervised fusion & A5 plus complete-cycle ROM loss\\
A7 & A6 + per-view-peak ROM anchor & A6 plus post hoc trunk-twist anchor\\
A8 & A7 + rotation-parameterized residual & A7 plus post hoc twist residual\\
A9 & A8 + observed twist-rate anchor & A8 plus post hoc rate anchor\\
\bottomrule
\end{tabular}
\end{table}

\section{Complete held-out learned comparison}

\begin{table}[h]
\caption{Held-out comparison ($N=14$). Mean per-joint position error is
reported after one similarity alignment per cycle followed by framewise hip
centring to the same-video triangulated pseudo-reference. Differences are
method minus A6 with participant-bootstrap 95\% confidence intervals and
Holm-adjusted paired Wilcoxon $p$ values. Rotation range of motion (ROM) and
peak angular speed are A3-relative retention ratios, not comparisons with
independent kinematics.}
\label{tab:complete-heldout}
\centering
\scriptsize
\setlength{\tabcolsep}{3pt}
\begin{tabular}{clccc}
\toprule
ID & Method & Error (mm) & Difference [95\% CI] &
$p_{\rm Holm}$\\
\midrule
A0 & Face only & $77.92\pm12.52$ & +17.14 [13.69, 21.34] & 0.0011\\
A1 & Side only & $75.84\pm13.93$ & +15.07 [10.23, 19.48] & 0.0018\\
A2 & Arithmetic mean & $60.85\pm8.98$ & +0.07 [$-0.50$, 0.59] & 1.0000\\
A3 & Quality mean & $61.03\pm8.84$ & +0.26 [0.17, 0.35] & 0.0011\\
A4 & Spatial objectives & $62.81\pm8.86$ & +2.04 [1.71, 2.36] & 0.0011\\
A5 & Rotation/temporal & $60.80\pm8.96$ & +0.02 [$-0.38$, 0.43] & 1.0000\\
A6 & Complete model & $60.78\pm8.72$ & -- & --\\
A7 & Per-view-peak ROM & $61.31\pm8.87$ & +0.54 [0.38, 0.69] & 0.0011\\
A8 & Twist residual & $92.02\pm9.59$ & +31.25 [27.78, 34.88] & 0.0011\\
A9 & Twist-rate anchor & $94.11\pm9.14$ & +33.34 [30.00, 36.77] & 0.0011\\
\bottomrule
\end{tabular}
\end{table}

All A2--A9 outputs were view-order invariant to the reported precision.
A3-relative (ROM, peak-speed) retention was A0 (1.156, 1.018), A1 (1.050,
1.654), A2 (1.101, 0.917), A3 (1.000, 1.000), A4 (0.949, 0.828), A5
(0.943, 0.644), A6 (1.000, 1.000), A7 (0.948, 0.821), A8 (1.218, 2.187) and
A9 (0.997, 2.615). The A8 and A9 variants demonstrate that relaxing the
conservative correction can increase coordinate error and over-amplify peak
angular speed.

\clearpage
\section{Deterministic comparison}

\begin{table}[h]
\caption{Agreement with the private triangulated pseudo-reference over 137
participants. Values are mm after one similarity alignment per cycle followed
by framewise hip centring. The pseudo-reference-fitted row is a leakage
diagnostic and is not an eligible label-free method.}
\label{tab:deterministic}
\centering
\scriptsize
\setlength{\tabcolsep}{3pt}
\begin{tabular}{p{0.42\linewidth}rrr}
\toprule
Method & Mean & SD & 95\% CI\\
\midrule
Body-frame average & 65.249 & 14.523 & [62.935, 67.789] \\ % deterministic-row
World-coordinate average & 123.703 & 18.096 & [120.710, 126.803] \\ % deterministic-row
Root alignment and average & 123.709 & 18.096 & [120.678, 126.787] \\ % deterministic-row
Similarity alignment, all joints & 69.317 & 18.250 & [66.358, 72.465] \\ % deterministic-row
Similarity alignment, stable joints & 66.086 & 14.291 & [63.719, 68.557] \\ % deterministic-row
Pseudo-reference-fitted joint weights (leaky) & 64.077 & 14.449 & [61.718, 66.539] \\ % deterministic-row
Similarity alignment, body-part weights & 66.270 & 14.593 & [63.863, 68.745] \\ % deterministic-row
Similarity alignment, side smoothing & 67.037 & 14.268 & [64.686, 69.437] \\ % deterministic-row
Similarity alignment, output smoothing & 68.031 & 14.298 & [65.731, 70.435] \\ % deterministic-row
\bottomrule
\end{tabular}
\end{table}


\input{generated/extrinsic_comparison_all137}

In this secondary all-participant analysis, Extrinsic-R reduced the mean by
2.175~mm relative to deterministic body-frame averaging and was lower for 109
of 137 participants ($p_{\rm Holm}=1.24\times10^{-14}$). The quality-weighted
variant reduced the mean by 1.116~mm and was lower for 89 participants
($p_{\rm Holm}=4.92\times10^{-5}$). Holdout reprojection error was associated
with Extrinsic-R error (Spearman $\rho=0.347$, $p=3.35\times10^{-5}$). These
results describe all available participants, including training and validation
participants, and are not an independent learned-model test.
\clearpage

\section{Synthetic corruptions and robustness}

\begin{table}[h]
\caption{Synthetic-corruption configuration in canonical units.}
\label{tab:corruptions}
\centering
\small
\begin{tabular}{ll}
\toprule
Corruption & Setting\\
\midrule
Independent joint dropout & probability 0.05\\
Temporal block dropout & joint probability 0.25; 16 frames\\
Impulse noise & probability 0.02; scale 0.10\\
Random-walk drift & joint probability 0.10; step scale 0.01\\
Thorax rotation bias & frame probability 0.10; $\pm12^\circ$\\
Frozen segment & joint probability 0.10; 12 frames\\
Integer time shift & joint probability 0.10; at most 4 frames\\
\bottomrule
\end{tabular}
\end{table}

\begin{table}[h]
\caption{Output displacement under the fixed corruption manifest (27
validation participants, one seed). Lower values indicate less movement from
each method's own clean output, not recovery toward independent ground truth.}
\label{tab:robustness}
\centering
\small
\begin{tabular}{lcc}
\toprule
Challenge & A3 & A6\\
\midrule
Joint dropout & 0.0629 & 0.0873\\
Temporal block dropout & 0.0619 & 0.0873\\
Impulse noise & 0.0605 & 0.0605\\
Random-walk drift & 0.0445 & 0.0445\\
Thorax rotation bias & 0.0312 & 0.0312\\
Frozen segment & 0.0643 & 0.0643\\
Integer time shift & 0.0348 & 0.0348\\
\bottomrule
\end{tabular}
\end{table}

A6 had greater aggregate clean-to-corrupted displacement than A4 and A5
(0.120 versus approximately 0.09 repository units). A temporal-offset sweep
over all 137 participants moved A6 by 0.046 and 0.078 units at offsets of
$\pm2$ and $\pm4$ frames, respectively, essentially matching the deterministic
mean. The learned model therefore did not resolve dependence on the upstream
global offset.

\section{Additional external comparisons}

Two Unity-supervised comparators were evaluated over two direction-held-out
folds and three seeds. An exact-extrinsic gate over the two 3D streams averaged
153.645~mm and 41.090$^\circ$, a 6.74\% improvement over its matched direct
average. Learnable triangulation averaged 28.058~mm and 14.555$^\circ$, 1.88\%
below its matched fixed-triangulation baseline. These rows use Unity native 3D
during training and answer a different question from zero-shot private-model
transfer.

The fitted-camera real-data pilot used the 14 held-out participants and three
seeds. The frozen A6 baseline averaged 60.578~mm; the full camera-feature model
averaged 60.850~mm, and its 30-degree wrong-camera control averaged
60.850~mm. Correct-camera features therefore failed the preregistered gate and
are retained as a negative result rather than a positive comparator.

\section{Complete per-joint comparison}

The following descriptive table uses the same 14 held-out participants and the
same sequence-similarity-plus-hip-centring evaluator for every method. It does
not constitute 70 separate hypothesis tests.

\begingroup
\scriptsize
\setlength{\tabcolsep}{2pt}
\begin{longtable}{lrrrrrr}
\caption{Complete MHR70 per-joint agreement on the same 14 held-out
participants. Values are participant-level mean MPJPE in mm; bold is the
lowest descriptive value in each row. Each cycle uses one similarity alignment
followed by framewise hip centring.}\label{tab:joint-accuracy-all70}\\
\toprule
Joint & Face & Side & Body mean & A6 & Extrinsic-R & Extrinsic-R quality \\
\midrule
\endfirsthead
\multicolumn{7}{c}{\tablename\ \thetable{} -- continued}\\
\toprule
Joint & Face & Side & Body mean & A6 & Extrinsic-R & Extrinsic-R quality \\
\midrule
\endhead
\midrule
\multicolumn{7}{r}{Continued on next page}\\
\endfoot
\bottomrule
\endlastfoot
nose & 65.3 & \textbf{38.8} & 40.8 & 40.0 & 39.9 & 41.3 \\ % joint-row
left-eye & 65.5 & 41.7 & 40.6 & \textbf{40.2} & 40.5 & 42.1 \\ % joint-row
right-eye & 64.0 & 41.7 & 40.2 & 39.3 & \textbf{39.2} & 41.1 \\ % joint-row
left-ear & 61.5 & 40.0 & 38.2 & \textbf{38.1} & 39.1 & 40.5 \\ % joint-row
right-ear & 56.9 & 39.2 & 36.7 & 35.6 & \textbf{35.4} & 37.0 \\ % joint-row
left-shoulder & 43.7 & 31.1 & 31.3 & \textbf{30.9} & 31.2 & 31.7 \\ % joint-row
right-shoulder & 40.0 & 29.0 & 29.7 & 27.9 & \textbf{27.0} & 27.5 \\ % joint-row
left-elbow & 47.1 & 47.9 & 38.4 & 39.1 & \textbf{36.8} & 38.1 \\ % joint-row
right-elbow & 47.8 & 47.8 & 38.2 & 37.9 & \textbf{35.4} & 35.8 \\ % joint-row
left-hip & 6.4 & 8.7 & 6.1 & 6.3 & \textbf{5.3} & 5.4 \\ % joint-row
right-hip & 6.4 & 8.7 & 6.1 & 6.3 & \textbf{5.3} & 5.4 \\ % joint-row
left-knee & 27.5 & 40.9 & 27.7 & 28.3 & \textbf{27.2} & 27.7 \\ % joint-row
right-knee & \textbf{21.6} & 36.1 & 25.7 & 24.8 & 24.4 & 24.4 \\ % joint-row
left-ankle & \textbf{42.5} & 73.0 & 53.0 & 51.5 & 51.6 & 51.7 \\ % joint-row
right-ankle & \textbf{37.0} & 65.4 & 49.6 & 47.6 & 47.9 & 47.4 \\ % joint-row
left-big-toe-tip & \textbf{47.1} & 82.7 & 57.2 & 56.7 & 57.2 & 57.6 \\ % joint-row
left-small-toe-tip & \textbf{45.1} & 81.4 & 55.9 & 55.5 & 55.4 & 55.8 \\ % joint-row
left-heel & \textbf{45.5} & 77.5 & 57.2 & 55.3 & 55.6 & 55.6 \\ % joint-row
right-big-toe-tip & \textbf{39.9} & 73.0 & 52.1 & 50.1 & 50.6 & 50.4 \\ % joint-row
right-small-toe-tip & \textbf{40.7} & 75.0 & 54.0 & 51.8 & 52.2 & 52.0 \\ % joint-row
right-heel & \textbf{40.7} & 69.4 & 53.5 & 51.4 & 51.8 & 51.3 \\ % joint-row
right-thumb-tip & 109.1 & 104.8 & 82.0 & 80.4 & 80.7 & \textbf{79.6} \\ % joint-row
right-thumb-first-joint & 104.8 & 99.8 & 78.7 & 77.1 & 77.3 & \textbf{76.3} \\ % joint-row
right-thumb-second-joint & 99.4 & 93.3 & 74.3 & 72.9 & 72.8 & \textbf{71.9} \\ % joint-row
right-thumb-third-joint & 93.8 & 87.4 & 70.0 & 68.8 & 68.3 & \textbf{67.6} \\ % joint-row
right-index-tip & 112.5 & 111.3 & 85.8 & 84.3 & 84.5 & \textbf{83.5} \\ % joint-row
right-index-first-joint & 110.2 & 108.7 & 84.1 & 82.6 & 82.7 & \textbf{81.7} \\ % joint-row
right-index-second-joint & 107.4 & 105.4 & 81.8 & 80.5 & 80.4 & \textbf{79.5} \\ % joint-row
right-index-third-joint & 102.2 & 98.6 & 77.4 & 76.1 & 75.9 & \textbf{75.1} \\ % joint-row
right-middle-tip & 110.1 & 109.6 & 83.9 & 82.6 & 82.4 & \textbf{81.5} \\ % joint-row
right-middle-first-joint & 108.0 & 107.3 & 82.4 & 81.2 & 80.9 & \textbf{80.1} \\ % joint-row
right-middle-second-joint & 106.0 & 104.9 & 80.9 & 79.7 & 79.3 & \textbf{78.6} \\ % joint-row
right-middle-third-joint & 101.2 & 98.1 & 76.5 & 75.6 & 74.8 & \textbf{74.2} \\ % joint-row
right-ring-tip & 107.6 & 107.1 & 81.7 & 80.5 & 80.0 & \textbf{79.2} \\ % joint-row
right-ring-first-joint & 105.7 & 104.9 & 80.3 & 79.2 & 78.6 & \textbf{77.9} \\ % joint-row
right-ring-second-joint & 103.8 & 102.3 & 78.8 & 77.9 & 77.0 & \textbf{76.5} \\ % joint-row
right-ring-third-joint & 99.1 & 96.0 & 74.7 & 73.9 & 72.8 & \textbf{72.4} \\ % joint-row
right-pinky-tip & 102.3 & 102.9 & 77.9 & 77.0 & 76.0 & \textbf{75.5} \\ % joint-row
right-pinky-first-joint & 101.0 & 101.0 & 76.8 & 76.0 & 74.9 & \textbf{74.4} \\ % joint-row
right-pinky-second-joint & 99.7 & 98.7 & 75.6 & 75.0 & 73.6 & \textbf{73.3} \\ % joint-row
right-pinky-third-joint & 96.7 & 93.7 & 72.7 & 72.1 & 70.7 & \textbf{70.4} \\ % joint-row
right-wrist & 89.3 & 83.2 & 66.4 & 65.6 & 64.4 & \textbf{64.1} \\ % joint-row
left-thumb-tip & 102.0 & 92.7 & 75.9 & 77.8 & \textbf{73.6} & 75.2 \\ % joint-row
left-thumb-first-joint & 98.1 & 88.1 & 72.7 & 74.6 & \textbf{70.5} & 72.1 \\ % joint-row
left-thumb-second-joint & 93.2 & 82.4 & 68.7 & 70.5 & \textbf{66.5} & 68.0 \\ % joint-row
left-thumb-third-joint & 88.2 & 77.0 & 64.8 & 66.5 & \textbf{62.5} & 64.1 \\ % joint-row
left-index-tip & 105.9 & 96.5 & 79.2 & 81.1 & \textbf{76.6} & 78.2 \\ % joint-row
left-index-first-joint & 104.0 & 94.7 & 77.9 & 79.8 & \textbf{75.2} & 76.9 \\ % joint-row
left-index-second-joint & 101.7 & 92.4 & 76.1 & 78.0 & \textbf{73.4} & 75.1 \\ % joint-row
left-index-third-joint & 96.8 & 86.6 & 71.9 & 73.7 & \textbf{69.3} & 70.9 \\ % joint-row
left-middle-tip & 105.2 & 96.2 & 78.5 & 80.5 & \textbf{75.9} & 77.6 \\ % joint-row
left-middle-first-joint & 103.4 & 94.2 & 77.3 & 79.3 & \textbf{74.6} & 76.3 \\ % joint-row
left-middle-second-joint & 101.5 & 92.0 & 75.8 & 77.8 & \textbf{73.1} & 74.8 \\ % joint-row
left-middle-third-joint & 96.5 & 85.6 & 71.5 & 73.3 & \textbf{68.8} & 70.4 \\ % joint-row
left-ring-tip & 103.1 & 94.1 & 76.8 & 78.8 & \textbf{74.2} & 75.9 \\ % joint-row
left-ring-first-joint & 101.3 & 92.0 & 75.5 & 77.5 & \textbf{72.9} & 74.6 \\ % joint-row
left-ring-second-joint & 99.4 & 89.5 & 74.1 & 75.9 & \textbf{71.4} & 73.1 \\ % joint-row
left-ring-third-joint & 94.9 & 83.6 & 70.1 & 71.9 & \textbf{67.5} & 69.1 \\ % joint-row
left-pinky-tip & 99.5 & 90.3 & 74.0 & 75.9 & \textbf{71.4} & 73.2 \\ % joint-row
left-pinky-first-joint & 98.1 & 88.5 & 73.0 & 74.9 & \textbf{70.4} & 72.1 \\ % joint-row
left-pinky-second-joint & 96.5 & 86.4 & 71.8 & 73.6 & \textbf{69.1} & 70.8 \\ % joint-row
left-pinky-third-joint & 93.0 & 81.5 & 68.6 & 70.3 & \textbf{66.0} & 67.6 \\ % joint-row
left-wrist & 84.8 & 72.9 & 62.0 & 63.7 & \textbf{59.6} & 61.2 \\ % joint-row
left-olecranon & 47.8 & 50.4 & 39.5 & 40.2 & \textbf{37.8} & 39.2 \\ % joint-row
right-olecranon & 47.8 & 50.1 & 39.0 & 38.9 & \textbf{36.2} & 36.9 \\ % joint-row
left-cubital-fossa & 47.9 & 47.1 & 38.4 & 39.0 & \textbf{36.8} & 38.1 \\ % joint-row
right-cubital-fossa & 49.6 & 48.1 & 38.9 & 38.4 & \textbf{36.2} & 36.4 \\ % joint-row
left-acromion & 47.0 & \textbf{31.8} & 32.8 & 32.6 & 32.9 & 33.5 \\ % joint-row
right-acromion & 42.4 & 29.1 & 30.6 & 28.7 & \textbf{28.0} & 28.6 \\ % joint-row
neck & 42.7 & \textbf{27.3} & 28.8 & 27.8 & 28.2 & 28.9 \\ % joint-row
\end{longtable}

\endgroup

\section{Complete cohort outcomes}

\begin{table}[h]
\caption{Out-of-fold cohort and repeated-cycle analysis. Cohort effects are
elderly-minus-student mixed-model coefficients at normalized cycle position
0.5. MAD effects are participant-level median differences.}
\label{tab:cohort}
\centering
\scriptsize
\setlength{\tabcolsep}{3pt}
\begin{tabular}{p{0.28\linewidth}p{0.30\linewidth}cc}
\toprule
Outcome & Cohort effect [95\% CI] & $p_{\rm Holm}$ &
MAD effect ($p_{\rm Holm}$)\\
\midrule
Axial rotation ROM & 0.2475 [$-0.0242$, 0.5192] & 0.2966 & $-0.0016$ (1.0000)\\
Log angular speed (P95) & 0.3006 [0.1136, 0.4876] & 0.0114 & 0.0314 (0.4459)\\
Peak rotation phase & $-0.0013$ [$-0.0568$, 0.0542] & 1.0000 & 0.0000 (1.0000)\\
Trunk tilt (P95) & 0.0042 [0.0002, 0.0082] & 0.2514 & $-0.0002$ (1.0000)\\
Wrist wrapping (P95) & 0.0289 [$-0.0626$, 0.1204] & 1.0000 & $-0.0073$ (1.0000)\\
Log cycle duration & $-0.0127$ [$-0.0256$, 0.0002] & 0.2677 & 0.0083 (1.0000)\\
Log dimensionless jerk & 0.4233 [0.2336, 0.6130] & 0.0001 & 0.0377 (0.0208)\\
Log whole-body repeatability & $-0.0145$ [$-0.1055$, 0.0765] & 1.0000 & $-0.0003$ (1.0000)\\
\bottomrule
\end{tabular}
\end{table}

All eight cohort-by-cycle-position interactions had
$p_{\rm Holm}=1.000$. No axial-angle or angular-speed phase cluster survived
correction across descriptor families. Secondary phase-localized differences
occurred for trunk tilt and wrist wrapping, but are hypothesis-generating and
not anatomical validation. Omitting outer-fold fixed effects produced similar
coefficients for log angular speed (0.2972) and log jerk (0.4161).

\end{document}
